\documentclass[aspectratio=169]{beamer}

\usetheme{Madrid}
\usecolortheme{default}

\usepackage{amsmath}
\usepackage{booktabs}
\usepackage{tikz}
\usepackage{graphicx}
\usepackage{hyperref}

\title{Modern Foundation Models}
\subtitle{Lecture 4: LLM Training, Optimization and Fine-Tuning}
\author{Dror Lederman}
\institute{Holon Institute of Technology}
\date{\today}

\begin{document}

%------------------------------------------------
\begin{frame}
\titlepage
\end{frame}

%------------------------------------------------
\begin{frame}{Lecture Roadmap}
\begin{enumerate}
    \item The LLM Training Paradigm
    \item Pretraining
    \item Scaling Laws and Compute
    \item Training Bottlenecks
    \item Distributed Training
    \item FlashAttention
    \item Mixed Precision Training
    \item Supervised Fine-Tuning
    \item Parameter-Efficient Fine-Tuning: LoRA and QLoRA
\end{enumerate}
\end{frame}

%------------------------------------------------
\section{LLM Training Paradigm}

\begin{frame}{Paradigm Shift}

\textbf{Traditional ML:}

\begin{itemize}
    \item Given a task, train a task-specific model.
    \item Example: one model for spam detection, another for sentiment, another for translation.
\end{itemize}

\vspace{0.5cm}

\textbf{Modern LLM paradigm:}

\begin{itemize}
    \item First train a general-purpose language model.
    \item Then adapt it to downstream tasks.
\end{itemize}

\end{frame}

%------------------------------------------------
\begin{frame}{Two-Stage LLM Training}

\[
\text{Raw Text + Code}
\rightarrow
\text{Pretraining}
\rightarrow
\text{Pretrained Model}
\rightarrow
\text{Fine-Tuning}
\rightarrow
\text{Task / Instruction Model}
\]

\vspace{0.5cm}

Main stages:

\begin{enumerate}
    \item \textbf{Pretraining:} learn general patterns of language and code.
    \item \textbf{Tuning:} adapt behavior to desired tasks.
\end{enumerate}

\end{frame}

%------------------------------------------------
\section{Pretraining}

\begin{frame}{Pretraining Goal}

The goal of pretraining is to learn statistical patterns of:

\begin{itemize}
    \item Natural language
    \item Code
    \item Multilingual text
    \item Structured and semi-structured data
\end{itemize}

\vspace{0.5cm}

A pretrained model is not trained for one specific task.

\vspace{0.3cm}

It learns a broad representation of text and code.

\end{frame}

%------------------------------------------------
\begin{frame}{Pretraining Objective}

For decoder-only LLMs, the main objective is:

\[
P(w_1,\ldots,w_T)
=
\prod_{t=1}^{T}
P(w_t \mid w_{<t})
\]

\vspace{0.5cm}

Equivalent training objective:

\[
\mathcal{L}
=
-
\sum_{t=1}^{T}
\log P(w_t \mid w_{<t})
\]

\vspace{0.4cm}

The model learns by predicting the next token.

\end{frame}

%------------------------------------------------
\begin{frame}{Next Token Prediction}

Example:

\[
\texttt{[BOS] A teddy bear is}
\rightarrow
?
\]

\vspace{0.5cm}

The model produces a probability distribution over the vocabulary:

\[
P(w_{t+1} \mid w_1,\ldots,w_t)
\]

\vspace{0.4cm}

The correct next token is used as the training target.

\end{frame}

%------------------------------------------------
\begin{frame}{Pretraining Data Mixtures}

Typical data sources:

\begin{itemize}
    \item Web text: Common Crawl, Wikipedia
    \item Books and articles
    \item Code: GitHub, StackOverflow
    \item Multilingual corpora
    \item Mathematical and scientific text
\end{itemize}

\vspace{0.5cm}

Modern LLMs are trained on hundreds of billions to trillions of tokens.

\end{frame}

%------------------------------------------------
\begin{frame}{Examples of Training Scale}

\begin{table}
\centering
\begin{tabular}{lc}
\toprule
Model & Approximate Pretraining Tokens \\
\midrule
GPT-3 & 300B \\
LLaMA 3 & 15T \\
\bottomrule
\end{tabular}
\end{table}

\vspace{0.5cm}

Training scale is determined by:

\begin{itemize}
    \item Number of parameters
    \item Number of tokens
    \item Available compute
\end{itemize}

\end{frame}

%------------------------------------------------
\section{Scaling Laws}

\begin{frame}{FLOPs and FLOPS}

\textbf{FLOPs}

\[
\text{FLoating Point OPerations}
\]

Total number of operations needed for computation.

\vspace{0.5cm}

\textbf{FLOPS}

\[
\text{FLoating Point OPerations per Second}
\]

Computational throughput of hardware.

\vspace{0.5cm}

Training LLMs requires enormous numbers of FLOPs.

\end{frame}

%------------------------------------------------
\begin{frame}{Scaling Laws}

Empirical scaling laws show that test loss often follows power-law behavior with respect to:

\begin{itemize}
    \item Model size
    \item Dataset size
    \item Training compute
\end{itemize}

\vspace{0.5cm}

General form:

\[
L(N,D,C)
\approx
A N^{-\alpha}
+
B D^{-\beta}
+
E C^{-\gamma}
\]

\vspace{0.3cm}

where $N$ is parameters, $D$ is data, and $C$ is compute.

\end{frame}

%------------------------------------------------
\begin{frame}{Sample Efficiency}

Larger models tend to learn more efficiently from data.

\vspace{0.5cm}

Observation:

\begin{itemize}
    \item Larger models often reach lower loss with fewer tokens.
    \item But they require more memory and compute.
    \item The optimal model size depends on the compute budget.
\end{itemize}

\end{frame}

%------------------------------------------------
\begin{frame}{Chinchilla Law}

The Chinchilla result argued that many earlier LLMs were undertrained.

\vspace{0.5cm}

Key message:

\begin{quote}
For a fixed compute budget, it is often better to train a smaller model on more tokens.
\end{quote}

\vspace{0.5cm}

Implication:

\[
\text{Model size and data size should scale together.}
\]

\end{frame}

%------------------------------------------------
\begin{frame}{Why Pretraining Is Difficult}

Main challenges:

\begin{itemize}
    \item Very high cost
    \item Long training time
    \item Massive GPU clusters
    \item Energy consumption
    \item Data quality and filtering
    \item Knowledge cutoff
    \item Difficult knowledge editing
\end{itemize}

\end{frame}

%------------------------------------------------
\section{Training Setup}

\begin{frame}{LLM Training Setup}

Training requires:

\begin{enumerate}
    \item Transformer model
    \item Large text/code dataset
    \item Many GPUs
    \item Distributed training framework
    \item Optimizer
    \item Checkpointing and monitoring
\end{enumerate}

\vspace{0.4cm}

The bottleneck is often not only compute, but also memory and communication.

\end{frame}

%------------------------------------------------
\begin{frame}{Training Step}

Each training step includes:

\begin{enumerate}
    \item Forward pass
    \item Loss computation
    \item Backward pass
    \item Gradient computation
    \item Optimizer update
\end{enumerate}

\vspace{0.5cm}

\[
\theta_{t+1}
=
\theta_t
-
\alpha \nabla_{\theta}\mathcal{L}(\theta_t)
\]

\end{frame}

%------------------------------------------------
\begin{frame}{Adam Optimizer}

Adam keeps moving averages of gradients and squared gradients:

\[
m_t
=
\beta_1 m_{t-1}
+
(1-\beta_1)g_t
\]

\[
v_t
=
\beta_2 v_{t-1}
+
(1-\beta_2)g_t^2
\]

\[
\theta_{t+1}
=
\theta_t
-
\alpha
\frac{m_t}{\sqrt{v_t}+\epsilon}
\]

\vspace{0.4cm}

Adam improves stability, but adds optimizer-state memory.

\end{frame}

%------------------------------------------------
\begin{frame}{Memory Bottleneck}

GPU memory must store:

\begin{itemize}
    \item Model parameters
    \item Activations
    \item Gradients
    \item Optimizer state
\end{itemize}

\vspace{0.5cm}

For large LLMs, memory becomes a central bottleneck.

\vspace{0.4cm}

Even modern GPUs have limited memory relative to LLM size.

\end{frame}

%------------------------------------------------
\section{Distributed Training}

\begin{frame}{Data Parallelism}

\textbf{Idea:}

\begin{itemize}
    \item Split the batch across GPUs.
    \item Replicate the full model on each GPU.
    \item Synchronize gradients across devices.
\end{itemize}

\vspace{0.5cm}

Pros:

\begin{itemize}
    \item Simple
    \item Effective for moderate-size models
\end{itemize}

Cons:

\begin{itemize}
    \item Replicates parameters, gradients and optimizer state.
\end{itemize}

\end{frame}

%------------------------------------------------
\begin{frame}{ZeRO: Zero Redundancy Optimization}

ZeRO reduces redundant memory across devices.

\vspace{0.5cm}

Stages:

\begin{itemize}
    \item \textbf{ZeRO-1:} partition optimizer state.
    \item \textbf{ZeRO-2:} partition optimizer state and gradients.
    \item \textbf{ZeRO-3:} partition optimizer state, gradients and parameters.
\end{itemize}

\vspace{0.5cm}

Goal: train larger models with the same GPU memory.

\end{frame}

%------------------------------------------------
\begin{frame}{Model Parallelism}

\textbf{Idea:}

Split the model computation across multiple GPUs.

\vspace{0.5cm}

Common forms:

\begin{itemize}
    \item Tensor Parallelism
    \item Pipeline Parallelism
    \item Sequence Parallelism
    \item Context Parallelism
    \item Expert Parallelism
\end{itemize}

\vspace{0.4cm}

Used when the model cannot fit on one GPU.

\end{frame}

%------------------------------------------------
\section{FlashAttention}

\begin{frame}{Why FlashAttention?}

Standard attention:

\[
\mathrm{Attention}(Q,K,V)
=
\mathrm{softmax}
\left(
\frac{QK^T}{\sqrt{d_k}}
\right)V
\]

\vspace{0.5cm}

Problem:

\begin{itemize}
    \item Attention creates large intermediate matrices.
    \item Moving data between GPU memory levels is expensive.
    \item HBM is large but slower; SRAM is small but fast.
\end{itemize}

\end{frame}

%------------------------------------------------
\begin{frame}{Standard Attention Computation}

A simplified standard computation:

\[
S = QK^T
\]

\[
P = \mathrm{softmax}(S)
\]

\[
O = PV
\]

\vspace{0.5cm}

Problem:

\begin{itemize}
    \item Store $S$ in high-bandwidth memory.
    \item Read $S$ again for softmax.
    \item Store $P$.
    \item Read $P$ again for multiplication with $V$.
\end{itemize}

\end{frame}

%------------------------------------------------
\begin{frame}{FlashAttention: Core Idea}

FlashAttention reduces memory reads and writes.

\vspace{0.5cm}

Main strategy:

\begin{quote}
Compute attention in blocks using fast SRAM, without materializing the full attention matrix in HBM.
\end{quote}

\vspace{0.5cm}

This is an exact attention algorithm, not an approximation.

\end{frame}

%------------------------------------------------
\begin{frame}{Tiling}

FlashAttention uses tiling:

\begin{itemize}
    \item Load blocks of $Q$, $K$, and $V$ into SRAM.
    \item Compute partial attention output.
    \item Accumulate result.
    \item Write only final output back to HBM.
\end{itemize}

\vspace{0.5cm}

Key benefit:

\[
\text{Less memory traffic}
\Rightarrow
\text{faster runtime}
\]

\end{frame}

%------------------------------------------------
\begin{frame}{Recompute Instead of Store}

In the backward pass, FlashAttention may recompute intermediate values.

\vspace{0.5cm}

Before:

\[
\text{Store } S,P \text{ in HBM}
\]

After:

\[
\text{Recompute } S,P \text{ using SRAM tiling}
\]

\vspace{0.5cm}

Important lesson:

\begin{quote}
More FLOPs can still mean less runtime if memory traffic is reduced.
\end{quote}

\end{frame}

%------------------------------------------------
\section{Precision and Mixed Precision}

\begin{frame}{Floating-Point Representation}

A floating-point number consists of:

\begin{itemize}
    \item Sign
    \item Exponent
    \item Mantissa
\end{itemize}

\vspace{0.5cm}

Different formats allocate different numbers of bits.

\end{frame}

%------------------------------------------------
\begin{frame}{Common Precision Formats}

\begin{table}
\centering
\begin{tabular}{lccc}
\toprule
Format & Sign & Exponent & Mantissa \\
\midrule
FP16 & 1 & 5 & 10 \\
FP32 & 1 & 8 & 23 \\
FP64 & 1 & 11 & 52 \\
BF16 & 1 & 8 & 7 \\
\bottomrule
\end{tabular}
\end{table}

\vspace{0.5cm}

Lower precision reduces memory and can accelerate computation.

\end{frame}

%------------------------------------------------
\begin{frame}{Mixed Precision Training}

Objective:

\begin{quote}
Speed up training and reduce memory requirements while preserving stability.
\end{quote}

\vspace{0.5cm}

Typical approach:

\begin{itemize}
    \item Forward pass in low precision.
    \item Backward pass in low precision.
    \item Keep master weights or critical updates in higher precision.
\end{itemize}

\end{frame}

%------------------------------------------------
\begin{frame}{Why Mixed Precision Works}

Benefits:

\begin{itemize}
    \item Less memory per tensor
    \item Faster matrix multiplication
    \item Larger batch sizes
    \item Better GPU utilization
\end{itemize}

\vspace{0.5cm}

Risk:

\begin{itemize}
    \item Numerical instability
    \item Underflow or overflow
\end{itemize}

\end{frame}

%------------------------------------------------
\section{Supervised Fine-Tuning}

\begin{frame}{Supervised Fine-Tuning}

\[
\text{SFT} = \text{Supervised Fine-Tuning}
\]

\vspace{0.5cm}

Goal:

\begin{quote}
Change the model behavior by tuning its weights on examples of desired behavior.
\end{quote}

\vspace{0.5cm}

SFT is usually performed after pretraining.

\end{frame}

%------------------------------------------------
\begin{frame}{SFT Data}

SFT data consists of input-output pairs:

\[
(x_i, y_i)
\]

\vspace{0.5cm}

Examples:

\begin{itemize}
    \item Instruction $\rightarrow$ answer
    \item Question $\rightarrow$ explanation
    \item Document $\rightarrow$ summary
    \item Prompt $\rightarrow$ desired structured output
\end{itemize}

\end{frame}

%------------------------------------------------
\begin{frame}{Instruction Tuning}

Instruction tuning is a special case of SFT.

\vspace{0.5cm}

The model is trained on instruction-following examples:

\[
\text{Instruction}
+
\text{Input}
\rightarrow
\text{Desired Response}
\]

\vspace{0.5cm}

This improves the model's ability to follow user intent.

\end{frame}

%------------------------------------------------
\begin{frame}{SFT Objective}

Given prompt $x$ and target response $y = (y_1,\ldots,y_T)$:

\[
\mathcal{L}_{SFT}
=
-
\sum_{t=1}^{T}
\log P(y_t \mid x, y_{<t})
\]

\vspace{0.5cm}

Same next-token prediction objective, but conditioned on task examples.

\end{frame}

%------------------------------------------------
\begin{frame}{Why Not Always Full Fine-Tuning?}

Full fine-tuning updates all model weights.

\vspace{0.5cm}

Problems:

\begin{itemize}
    \item Expensive
    \item Requires large GPUs
    \item Stores full optimizer state
    \item May overfit or forget prior capabilities
    \item Hard to maintain many task-specific variants
\end{itemize}

\end{frame}

%------------------------------------------------
\section{LoRA}

\begin{frame}{Parameter-Efficient Fine-Tuning}

Parameter-efficient fine-tuning updates only a small number of parameters.

\vspace{0.5cm}

Motivation:

\begin{itemize}
    \item Reduce memory
    \item Reduce compute
    \item Enable fine-tuning on smaller hardware
    \item Maintain a frozen base model
\end{itemize}

\vspace{0.5cm}

Popular method: LoRA.

\end{frame}

%------------------------------------------------
\begin{frame}{LoRA: Low-Rank Adaptation}

LoRA represents the weight update as a low-rank matrix:

\[
W
=
W_0
+
\Delta W
\]

\[
\Delta W
=
BA
\]

where:

\[
B \in \mathbb{R}^{d \times r},
\quad
A \in \mathbb{R}^{r \times k},
\quad
r \ll \min(d,k)
\]

\end{frame}

%------------------------------------------------
\begin{frame}{LoRA Training}

During LoRA fine-tuning:

\begin{itemize}
    \item Freeze the original weights $W_0$.
    \item Train only $A$ and $B$.
    \item Add the low-rank update during the forward pass.
\end{itemize}

\vspace{0.5cm}

\[
h = W_0x + BAx
\]

\vspace{0.5cm}

This greatly reduces trainable parameters.

\end{frame}

%------------------------------------------------
\begin{frame}{Why Low Rank?}

The assumption:

\begin{quote}
Task-specific adaptation often lies in a low-dimensional subspace.
\end{quote}

\vspace{0.5cm}

Therefore, a small-rank update may be enough to adapt the model.

\vspace{0.5cm}

Typical ranks:

\[
r = 4, 8, 16, 32, 64
\]

\end{frame}

%------------------------------------------------
\begin{frame}{Where to Apply LoRA?}

Common locations:

\begin{itemize}
    \item Query projection
    \item Key projection
    \item Value projection
    \item Output projection
    \item Feed-forward layers
\end{itemize}

\vspace{0.5cm}

In practice, applying LoRA to attention projections is very common.

\end{frame}

%------------------------------------------------
\begin{frame}{LoRA Tradeoffs}

Advantages:

\begin{itemize}
    \item Much fewer trainable parameters
    \item Lower memory footprint
    \item Easy to swap adapters
    \item Practical for many downstream tasks
\end{itemize}

\vspace{0.4cm}

Limitations:

\begin{itemize}
    \item Hyperparameter-sensitive
    \item May underperform full fine-tuning on some tasks
    \item Requires careful learning-rate tuning
\end{itemize}

\end{frame}

%------------------------------------------------
\section{QLoRA}

\begin{frame}{QLoRA}

QLoRA combines:

\begin{itemize}
    \item Quantized frozen base model
    \item Full-precision LoRA adapters
\end{itemize}

\vspace{0.5cm}

Idea:

\begin{quote}
Store the frozen base weights in low precision, but train LoRA adapters normally.
\end{quote}

\end{frame}

%------------------------------------------------
\begin{frame}{QLoRA Computation}

The model uses:

\[
W_0^{quantized}
+
BA
\]

\vspace{0.5cm}

where:

\begin{itemize}
    \item $W_0$ is stored quantized.
    \item $A$ and $B$ are stored and trained in higher precision.
    \item Computation may be performed in higher precision.
\end{itemize}

\vspace{0.5cm}

This greatly reduces memory requirements.

\end{frame}

%------------------------------------------------
\begin{frame}{Efficient Quantization}

Quantization maps floating-point weights to fewer bits.

\vspace{0.5cm}

Examples:

\begin{itemize}
    \item 8-bit quantization
    \item 4-bit quantization
    \item NF4: NormalFloat 4-bit
\end{itemize}

\vspace{0.5cm}

NF4 is designed for values that roughly follow a normal distribution.

\end{frame}

%------------------------------------------------
\begin{frame}{LoRA vs QLoRA}

\begin{table}
\centering
\begin{tabular}{lll}
\toprule
Method & Base Model & Trainable Parameters \\
\midrule
Full Fine-Tuning & Full precision & All weights \\
LoRA & Frozen full precision & Low-rank adapters \\
QLoRA & Frozen quantized & Low-rank adapters \\
\bottomrule
\end{tabular}
\end{table}

\vspace{0.5cm}

QLoRA is especially useful when GPU memory is limited.

\end{frame}

%------------------------------------------------
\begin{frame}{Key Takeaways}

\begin{itemize}
    \item LLM training usually follows pretraining plus tuning.
    \item Pretraining learns next-token prediction over massive text/code corpora.
    \item Scaling laws connect model size, data size and compute.
    \item Training is often memory-bound, not only compute-bound.
    \item ZeRO and model parallelism enable distributed training.
    \item FlashAttention reduces memory traffic.
    \item Mixed precision reduces memory and accelerates training.
    \item SFT changes model behavior using desired examples.
    \item LoRA and QLoRA enable practical fine-tuning of large models.
\end{itemize}

\end{frame}

%------------------------------------------------
\begin{frame}{Suggested Reading}

\begin{itemize}
    \item Brown et al. (2020), \textit{Language Models are Few-Shot Learners}
    \item Kaplan et al. (2020), \textit{Scaling Laws for Neural Language Models}
    \item Hoffmann et al. (2022), \textit{Training Compute-Optimal Large Language Models}
    \item Rajbhandari et al. (2019), \textit{ZeRO: Memory Optimizations Toward Training Trillion Parameter Models}
    \item Dao et al. (2022), \textit{FlashAttention}
    \item Micikevicius et al. (2017), \textit{Mixed Precision Training}
    \item Hu et al. (2021), \textit{LoRA}
    \item Dettmers et al. (2023), \textit{QLoRA}
\end{itemize}

\end{frame}

%------------------------------------------------
\begin{frame}
\centering
\Huge Questions?
\end{frame}

\end{document}